\documentclass[french]{book}
\usepackage{teach}
\usepackage{verbatim}
\usepackage{amsmath,amsfonts,amssymb}
\usepackage{pgfplots}
\pgfplotsset{compat=newest}
%\usepackage{geometry}
%\geometry{top=1cm, bottom=1cm, left=2cm, right=2cm}
\usepackage[utf8]{inputenc}
\usepackage[french]{babel}
\redefinecolor{thm}{title}{bgcolor}{alizarin}
\redefinecolor{prop}{title}{bgcolor}{meatbrown}
\redefinecolor{defn}{title}{bgcolor}{olivine}
\definecolor{alizarin}{rgb}{0.82, 0.1, 0.26}
\definecolor{meatbrown}{rgb}{0.9, 0.72, 0.23}
\definecolor{olivine}{rgb}{0.6, 0.73, 0.45}
\usepackage{mathrsfs}
%%
\newcommand{\R}{\mathbb {R}}
%\newcommand{\C}{\mathds {C}}
\newcommand{\N}{\mathbb {N}}
\newcommand{\Z}{\mathbb {Z}}
\newcommand{\Q}{\mathbb {Q}}
\newcommand{\D}{\mathbb {D}}
%%
\newcommand{\se}{\geq}
\newcommand{\ie}{\leq}
\newcommand{\tm}{\times}
%\newcommand{\ell}{l}
%\newcommand{\ell'}{l'}
%%
\newcommand{\limn}{\lim_{n\rightarrow +\infty}}
\newcommand{\limo}[1][x]{\ds\lim_{#1\rightarrow 0}}
\newcommand{\limite}[2][x]{\ds\lim_{#1\rightarrow #2}}
\newcommand{\limpinf}[1][x]{\ds\lim_{#1\rightarrow +\infty}}
\newcommand{\limminf}[1][x]{\ds\lim_{#1\rightarrow -\infty}}
\newcommand{\limiteg}[2][x]{\ds\lim_{#1\xrightarrow{<}#2}}
\newcommand{\limited}[2][x]{\ds\lim_{#1\xrightarrow{>}#2}}
%%
\newcommand{\suite}[1][u]{\left( #1_{n} \right)_{n\in\N}}
\newcommand{\suitep}[1][u]{\left( #1_{n} \right)_{n\in\N^{*}}}
\usetikzlibrary{arrows}


\begin{document}
\tableofcontents
\chapter{Suites numériques}

\section{Généralités}
\subsection{Définition}
\begin{defn}
Une suite est une fonction définie sur $\N$. Une suite numérique est une suite à valeurs dans $\R$.

 L'image de $n$ par une suite $u$ se note $u_{n}$ et est appelé terme de rang $n$ de la suite. Une suite $u$ est aussi notée $\suite[u]$.\\
 
Autrement dit, 
\begin{center}
$\begin{array}{ccccc}
u & : & \mathbb{N} & \to & \mathbb{R} \\
 & & n & \mapsto & u_n\\
\end{array}$
\end{center}
\end{defn}

\begin{rem}
On peut aussi définir une suite à partir d'un certain rang.
\end{rem}

\begin{exemple}
\begin{enumerate}
\item $u$ définie sur $\N$ par $u_{n}=\dfrac{1}{n+2}$.
\hspace{2em} $u_0=\dfrac{1}{2}$, $u_1=\dfrac{1}{3}$, $u_2=\dfrac{1}{4}$, ...
\item $v$ définie pour tout $n\se3$ par $v_{n}=\dfrac{1}{n-2}$
\hspace{2em} $v_3=1$, $v_4=\dfrac{1}{2}$, $v_5=\dfrac{1}{3}$, ...
\end{enumerate}
\end{exemple}


\subsection{Comment générer une suite}
Selon le contexte les termes d'une suite peuvent être définis de différentes façons.
\subsubsection*{Explicitement en fonction du rang}
\begin{itemize}


\item Toute fonction définie sur $[0~;~+\infty[$ (ou sur un intervalle de la forme $[a~;~+\infty[$) permet de définir une suite.
\end{itemize}
\begin{exemple}
Calculer les premiers termes de la suite définie sur $\N$ par $u_{n}=2n^2-5n-1$.
\begin{itemize}
\item $u_0=2\times0^2-5\times0-1=-1$ ; $u_1=2\times1^2-5\times1-1=-4$ ;
\item $u_2=2\times2^2-5\times2-1=-3$ ; $u_3=2\times3^2-5\times3-1=2$ ;
\item $u_4=2\times4^2-5\times4-1=11$
\end{itemize}
\end{exemple}

\begin{rem} 
Les propriétés des nombres entiers permettent aussi de définir explicitement des suites qui ne peuvent pas être obtenues simplement par une fonction définie sur $\R$.
\end{rem}
\begin{exemple}
Calculer les premiers termes des suites $\suite[u]$ et $\suitep[v]$ où $u_{n}=(-1)^{n}$ et $v_{n}$ est le nombre de diviseurs de $n$.

\begin{itemize}

\item $u_0=1$ ; $u_1=-1$ ; $u_2=1$ ; $u_3=-1$ ; ...
\vspace*{1.5ex}
\item $v_1=1$ ; $v_2=2$ ; $v_3=2$ ; $v_4=3$ ; $v_5=2$ ; $v_6=4$ ...
\end{itemize}
\end{exemple}

\subsubsection*{Par récurrence}
Une suite peut aussi être définie par son premier terme (ou ses premiers termes) et par une relation permettant de calculer chaque terme en fonction du précédent (ou des précédents).

\begin{exemple}
Calculer les premiers termes des suites ci-dessous définies sur $\N$.\\

$u$ est définie par
$
\left\{\begin{array}{l}
u_0=-6\\
u_{n+1}=-\dfrac{1}{2}u_{n}-1 \quad \text{pour tout }n\se0
\end{array}
\right.
$

\begin{itemize}
\item $u_1=-\dfrac{1}{2}u_0-1=-\dfrac{1}{2}\times(-6)-1=2$ ; $u_2=-\dfrac{1}{2}u_1-1=-\dfrac{1}{2}\times2-1=-2$
\item $u_3=-\dfrac{1}{2}u_2-1=-\dfrac{1}{2}\times(-2)-1=0$ ; $u_4=-\dfrac{1}{2}u_3-1=-\dfrac{1}{2}\times0-1=-1$
\end{itemize}

$v$ est définie par
$
\left\{\begin{array}{l}
v_0=1\; ; \; v_1=1\\
v_{n+2}=v_{n+1}+v_{n} \quad \text{pour tout }n\se0
\end{array}
\right.
$

\begin{itemize}
\item $v_2=v_1+v_0=1+1=2$ ; $v_3=v_2+v_1=2+1=3$ ;
\item $v_4=v_3+v_2=3+2=5$ ; $v_5=v_4+v_3=5+3=8$ ;
\item  $v_6=v_5+v_4=8+5=13$
\end{itemize}

$w$ est définie par
$w_0=3
\quad
\text{et}
\quad
\begin{cases}
w_{n+1}=\dfrac{w_{n}}{2}&\text{si $w_n$ est pair}\\
w_{n+1}=3w_{n}+1&\text{si $w_n$ est impair}
\end{cases}
$


\begin{itemize}
\item $w_1=3\times w_0 + 1 = 3\times3 +1 = 10$ ; $w_2=\dfrac{w_1}{2}=\dfrac{10}{2}=5$ ; \item $w_3=3\times w_2 + 1 = 3\times5 +1 = 16$ ; $w_4=\dfrac{w_3}{2}=\dfrac{16}{2}=8$ ; \item $w_5=\dfrac{w_4}{2}=\dfrac{8}{2}=4$ ; $w_6=\dfrac{w_5}{2}=\dfrac{4}{2}=2$
\end{itemize}
\end{exemple}
\newpage
\section{Sens de variation}
\subsection{Définition}
Une suite étant une fonction, les définitions restent les mêmes :

%\begin{tabular}{||l}
%$u$ est croissante (strict. croissante) si $n<p \Longrightarrow u_{n}\ie u_{p} \;(n<p \Longrightarrow u_{n}< u_{p})$\\
%$u$ est décroissante (strict. décroissante) si $n<p \Longrightarrow u_{n}\se u_{p} \;(n<p \Longrightarrow u_{n}> u_{p})$
%\end{tabular}

Cependant, les propriétés des nombres entiers permettent d'établir les résultats suivants :

\begin{prop}
$u$ est croissante si et seulement si pour tout entier $n$, $u_{n}\ie u_{n+1}$
\par $u$ est décroissante si et seulement si pour tout entier $n$, $u_{n}\se u_{n+1}$
\end{prop}

\begin{rem}
\begin{itemize}
\item Ne pas mélanger $u_{n+1}$ et $u_{n}+1$
\item Dans la pratique, pour déterminer le sens de variation d'une suite, on s'intéresse au signe de la différence $u_{n+1}-u_{n}$.
\item Dans le cas où tous les termes de la suite sont strictement positifs, on peut aussi comparer le quotient $\dfrac{u_{n+1}}{u_{n}}$ à 1.
\end{itemize}
\end{rem}

\begin{exemple}
Déterminer le sens de variation des suites suivantes :
 $u$ est définie sur $\N$ par $u_{n}=3n+(-1)^{n}$.
 Étudions la différence $u_{n+1}-u_{n}$ :
 $$\forall n\in\N, u_{n+1}-u_{n}=3(n+1)+(-1)^{n+1}-3n-(-1)^{n}=3-2\times(-1)^{n}>0$$
 La suite $\suite[u]$ est donc croissante.
\end{exemple}

\begin{exemple}
$v$ est définie sur $\N$ par 
$\left\{\begin{array}{l}
v_0=2\\
v_{n+1}=-v_{n}^2+v_{n} \quad \text{pour tout } n\in\N
\end{array}\right.$

 Étudions la différence $v_{n+1}-v_{n}$ :
 Pour tout $n\in\N$, $v_{n+1}-v_{n}=-v_{n}^2+v_{n}-v_{n}=-v_{n}^2<0$
 La suite $\suite{v}$ est donc décroissante.
\end{exemple}

\begin{exemple}
$w$ est définie sur $\N^{*}$ par $w_{n}=\dfrac{2^{n}}{n}$.


Les termes de la suite sont strictement positifs, on étudie donc le quotient $\dfrac{w_{n+1}}{w_{n}}$.
 Pour tout $n\in\N^{*}$, $\dfrac{w_{n+1}}{w_{n}} = \dfrac{\dfrac{2^{n+1}}{n+1}}{\dfrac{2^{n}}{n}} = \dfrac{2^{n+1}}{2^{n}}\times\dfrac{n}{n+1} = \dfrac{2n}{n+1}$
 Pour tout $n\se 1$, $2n\se n+1$ donc $\dfrac{w_{n+1}}{w_{n}}\se 1$.
 La suite $\suitep{w}$ est donc décroissante.

\end{exemple}


\subsection{Propriété}
\begin{prop}
Soit $f$ une fonction définie sur $\mathbb{R}_+$ et $u$ la suite définie sur $\N$ par $u_{n}=f(n)$.
\begin{itemize}


\item Si $f$ est croissante sur $\mathbb{R}_+$ alors $u$ est croissante.
\item Si $f$ est décroissante sur $\mathbb{R}_+$ alors $u$ est décroissante.
\end{itemize}
\end{prop}


\section{Limites}
\subsection{Suites convergentes}
\begin{defn}
Une suite $u$ converge vers un réel $\ell$ si tout intervalle ouvert contenant $\ell$ contient aussi tous les termes de la suite à partir d'un certain rang.
\par On dit alors que $u$ est convergente et on note $\limn u_{n}=l$
\end{defn}

On peut formuler la définition comme suit :
\begin{itemize}
\item $u$ converge vers $\ell$ si tout intervalle ouvert contenant $\ell$ contient tous les termes de la suite sauf un nombre fini.
\item $u$ converge vers $\ell$ si pour tout intervalle ouvert $I$ contenant $\ell$, il existe un entier $n_{0}$ tel que $n\se n_0\Longrightarrow u_n\in I$.
\end{itemize}

\vspace{2ex}
\begin{prop}
Si une suite converge alors sa limite est unique.
\end{prop}

\begin{demo}
Soit $u$ une suite convergente. Supposons qu'elle admette deux limites distinctes $\ell$ et $\ell'$ avec $\ell<\ell'$ sans perdre de généralité.\\

Soit $d$ un réel positif inférieur à $\dfrac{\ell'-\ell}{2}$.

On pose $I=]\ell-d~;~\ell+d[$ et $I'=]\ell'-d~;~\ell'+d[$.

On a $d<\dfrac{\ell'-\ell}{2}$ donc $2d<\ell'-\ell$ soit $d+d<\ell'-\ell$.

On en déduit que $\ell+d<\ell'-d$.\\

Ainsi, $I$ et $I'$ sont deux intervalles disjoints et ils contiennent respectivement $\ell$ et $\ell'$.

Si $u$ converge vers $\ell$, l'intervalle $I$ contient tous les termes de la suite à partir d'un rang $n_{0}$.

Si $u$ converge vers $\ell'$, l'intervalle $I'$ contient tous les termes de la suite à partir d'un rang $n_{1}$.

Ainsi pour tout $n$ supérieur à $n_{0}$ et $n_{1}$, $u_{n}$ doit appartenir à la fois à $I$ et $I'$ ce qui est impossible car ces intervalles sont disjoints. \\

\underline{ $u$ ne peut pas admettre deux limites distinctes. La limite d'une suite est unique.}
\end{demo}
\newpage

\subsection{Suites divergentes}
Une suite divergente est une suite qui n'est pas convergente.
\subsubsection*{Suites de limite infinie}
\begin{defn}
Une suite $u$ a pour limite $+\infty$ si tout intervalle ouvert de la forme $]A~;~+\infty[$ contient tous les termes de la suite à partir d'un certain rang.
\par On dit alors que $u$ diverge vers $+\infty$ et on note $\limn u_{n}=+\infty$
\par On définit de la même façon une suite qui diverge vers $-\infty$
\end{defn}
\subsubsection*{Suites qui n'ont pas de limite}
Une suite n'est pas nécessairement convergente ou divergente vers l'infini. Il existe des suites qui n'ont pas de limite. Elles sont aussi appelées suites divergentes.

\begin{exemple}
% \begin{itemize}
% \item 
$u$ définie sur $\N$ par $u_{n}=(-1)^{n}$.

%\item 
$v$ définie sur $\N$ par $v_{n}=n\sin(n)$.
%\end{itemize}
\end{exemple}

%\subsection{Propriétés}
%\begin{prop}
%Soit $f$ une fonction définie sur $[0~;~+\infty[$ et soit $u$ la suite définie sur $\N$ par $u_{n}=f(n)$.
%\par Si $\limpinf{x}f(x)=\ell$ alors $\limn u_{n}=\ell$.
%\par Si $\limpinf{x}f(x)=+\infty$ alors $\limn u_{n}=+\infty$.
%\end{prop}

%Conséquence : Toutes les propriétés vues sur les limites de fonctions s'appliquent aux suites définies par $u_{n}=f(n)$.% En particulier les limites usuelles, les règles de calcul sur la somme, le produit...
\begin{comment}
\begin{ex}
Déterminer la limite de la suite $u$ définie sur $\N$ par $u_{n}=\dfrac{-2n^{2}+n}{3n^{2}+1}$.
\end{ex}

\begin{repex}
Pour tout $n>0$, $u_{n}=\dfrac{-2n^{2}+n}{3n^{2}+1}=\dfrac{n^2\left( -2+\dfrac{1}{n} \right)}{n^2\left( 3+\dfrac{1}{n^2} \right)}=\dfrac{ -2+\dfrac{1}{n}}{3+\dfrac{1}{n^2}}$
\par $\left.\begin{array}{l}
\limn -2+\dfrac{1}{n}=-2\\
\limn 3+\dfrac{1}{n^2}=3
\end{array}\right\}$
donc $\limn u_{n}=-\dfrac{2}{3}$
\end{repex}
\end{comment}

\subsection{Théorème des gendarmes}

\begin{prop}
On considère trois suites $u$, $v$ et $w$ et un réel $\ell$.
\par Si \quad
$\left\{\begin{array}{l}
\limn u_{n}=\limn w_{n}=\ell \\
u_{n} \ie v_{n} \ie w_{n} \quad \text{à partir d'un certain rang}
\end{array}\right.
$
\qquad alors \quad $\limn v_{n}=\ell$.
\end{prop}

\begin{demo}
Soit $I$ un intervalle ouvert contenant $\ell$.
\par $u$ converge vers $\ell$ donc $I$ contient tous ses termes à partir d'un certain rang $n_{0}$.
\par $w$ converge vers $\ell$ donc $I$ contient tous ses termes à partir d'un certain rang $n_{1}$.
\par À partir d'un certain rang $n_{2}$, $u_{n} \ie v_{n} \ie w_{n}$.
\par Ainsi, pour tout $n$ supérieur à la fois à $n_{0}$, $n_{1}$ et $n_{2}$, tous les termes de $v$ appartiennent à $I$.
\par Conclusion : $v$ converge vers $\ell$.
\end{demo}

\newpage
\begin{exemple}
Déterminer la limite de la suite $u$ définie sur $\N^{*}$ par $u_{n}=2+\dfrac{(-1)^{n}}{n}$
\end{exemple}

\begin{align*}
\text{Pour tout } n \in \N \quad
&-1 \ie (-1)^n \ie 1 \\
&-\dfrac{1}{n} \ie \dfrac{(-1)^{n}}{n} \ie \dfrac{1}{n}\\
&2-\dfrac{1}{n} \ie 2+\dfrac{(-1)^{n}}{n} \ie 2+\dfrac{1}{n}\\
&2-\dfrac{1}{n} \ie u_n \ie 2+\dfrac{1}{n}
\end{align*}
$\left.\begin{array}{l}
\limn 2-\dfrac{1}{n}=2\\
\limn 2+\dfrac{1}{n}=2
\end{array}\right\}$
donc, d'après le théorème des gendarmes, $\limn u_n=2$


\section{Suites arithmétiques}
\subsection{Définition}
\begin{defn}
Soit $\suite[u]$ une suite et $r$ un réel.
 
On dit qu'une suite $\suite[u]$ est une suite arithmétique de raison $r$ si pour tout entier naturel $n$ : \[u_{n+1}=u_n+r\]
\end{defn}

\begin{exemple}
La suite des nombres impairs est une suite arithmétique de raison 2.

Les suites $u$ et $v$ définies sur $\N$ par $u_n=3n-4$ et $v_n=n^2+2$ sont-elles arithmétiques ?

\begin{itemize}
\item Pour tout $n$, $u_{n+1}-u_n=3(n+1)-4-(3n-4)=3n+3-4-3n+4=3$.

La suite $u$ est donc arithmétique de raison 3.

\item $v_0=2$, $v_1=3$, $v_2=6$ ainsi $v_1=v_0+1$ et $v_2=v_1+3$

La suite $v$ n'est donc pas arithmétique.
\end{itemize}
\end{exemple}

\vspace{2ex}
\begin{prop}
Soit $\suite[u]$ une suite arithmétique de raison $r$.
\par Si $r>0$ alors $u$ est strictement croissante.
\par Si $r=0$ alors $u$ est constante.
\par Si $r<0$ alors $u$ est strictement décroissante.
\end{prop}

\begin{demo}
Pour tout $n$, $u_{n+1}-u_n=r$
donc si  $r>0$ alors $u$ est strictement croissante, si $r=0$ alors $u$ est constante et si $r<0$ alors $u$ est strictement décroissante.
\end{demo}
\subsection{Expression de $u_n$ en fonction de $n$}
\begin{prop}
Soit $\suite[u]$ une suite arithmétique de raison $r$.
\par Pour tous entiers naturels $n$ et $p$ : \[u_n=u_p+(n-p)r\]
\par En particulier, pour tout entier naturel $n$ : \[u_n=u_0+nr\]
\end{prop}

\begin{demo}
On suppose $n>p$. On a : 
$u_{n}-u_{n-1}=r$ ; \hfill $u_{n-1}-u_{n-2}=r$ ; \hfill \dots \hfill $u_{p+2}-u_{p+1}=r$ ; \hfill $u_{p+1}-u_{p}=r$ \\
En additionnant toutes ces égalités, on obtient :
\[u_{n}-u_{p}=(n-p)r\]
soit \[u_n=u_p+(n-p)r\]
\end{demo}

\subsection{Somme de termes consécutifs}
\begin{prop}
Soit $\suite[u]$ une suite arithmétique de raison $r$.
\par Pour tout entier naturel $n$ :
\[u_{0}+u_{1}+ \cdots + u_{n-1}+u_{n}=(n+1)\times \dfrac{u_{0}+u_{n}}{2}\]
\end{prop}

\begin{rem}
On a aussi $u_{1}+u_2 \cdots + u_{n-1}+u_{n}=n\times \dfrac{u_{1}+u_{n}}{2}$
\end{rem}

\begin{demo}
\par Posons $S=u_0+u_1+ \cdots +u_n$
\par On a ainsi $2S=(u_0+u_n)+(u_1+u_{n-1})+(u_2+u_{n-2})+ \cdots +(u_{n-1}+u_1)+(u_n+u_0)$
\par Or, pour tout $k\ie n$, $u_k+u_{n-k}=u_0+kr+u_0+(n-k)r=u_0+u_0+nr=u_0+u_n$.
\par On a donc : $2S=(n+1)\times (u_0+u_n)$
\par Donc $S=(n+1)\times \dfrac{u_{0}+u_{n}}{2}$
\end{demo}

\newpage

\begin{exemple}
\underline{Déterminer la somme des $n$ premiers entiers naturels.}

La suite des entiers naturels est la suite arithmétique $\suitep[g]$ de raison 1 telle que $g_1=1$
\begin{equation*}
g_{1}+g_2 \cdots + g_{n-1}+g_{n}=n\times \dfrac{g_{1}+g_{n}}{2}=n\times \dfrac{g_{1}+g_{1}+(n-1)r}{2}=n\times\dfrac{1+(1+(n-1))}{2}=\dfrac{n(n+1)}{2}
\end{equation*}

\underline{Déterminer la somme des $n$ premiers nombres pairs.}

La suite des nombres pairs est la suite arithmétique $\suite[p]$ de raison 2 telle que $p_0=0$
\begin{equation*}
p_{0}+p_1 \cdots + p_{n-1}+p_{n}=n\times \dfrac{p_{1}+p_{n}}{2}=n\times \dfrac{p_{1}+p_{1}+(n+1)r}{2}=n\times\dfrac{0+0+(n+1)2}{2}=n(n+1)
\end{equation*}


\underline{Déterminer la somme des $n$ premiers nombres impairs.}

La suite des nombres impairs est la suite arithmétique $\suitep[i]$ de raison 2 telle que $i_1=1$
\begin{equation*}
i_{1}+i_2 \cdots + i_{n-1}+i_{n}=n\times \dfrac{i_{1}+i_{n}}{2}=n\times \dfrac{i_{1}+i_{1}+(n-1)r}{2}=n\times\dfrac{2+2n-2}{2}=n^2
\end{equation*}
\end{exemple}

\newpage

\subsection*{Application}
Pour finir cette exemple, voici une application de ces résultats, avec deux formules élégantes. 

 \paragraph*{Soit $n \in \N$ pair}
 
$n$ est égal à $p_{\frac{n}{2}}=n$, donc le  $\dfrac{n}{2}$ entier pair par définition de $p$.
\par Par suite il vient que $n-1$ est un entier impair, il est égal à $i_{\frac{n}{2}}=n$, donc le  $\dfrac{n}{2}$ entier impair par définition de $i$.

Ainsi on a les équations suivantes:
$$1+2+3+\cdots+(n-1)+n$$
$$\Longleftrightarrow i_1+ p_1+ i_2 + \cdots + i_{\frac{n}{2}} + p_{\frac{n}{2}}$$
quitte à réarranger les termes de la somme on a:
$$\Longleftrightarrow (p_1+p_2+\cdots+p_{\frac{n}{2}}) + (i_1+ i_2 + \cdots + i_{\frac{n}{2}})$$
$$\Longleftrightarrow (\dfrac{n}{2}(\dfrac{n}{2}+1))+ ( (\dfrac{n}{2})^2)$$

d'autre part: 
$$1+2+3+\cdots+(n-1)+n=\dfrac{n(n+1)}{2}$$
d'où 
$$\dfrac{n(n+1)}{2} =\dfrac{n}{2}(\dfrac{n}{2}+1) +  (\dfrac{n}{2})^2 $$
Cette formule nous dit simplement que la somme des entiers naturel $1+2+3+\cdots+n$ est égale à la somme de tous les entiers pairs et de tous les entiers impairs contenu dans $\{1;2;3;...;n\}$.


Avec nos suites on peut dire que $$g_{2n}=p_{n}+i_{n}$$

\paragraph*{Soit $n \in \N$ impair}

Si $n$ est impair, il est égal à $i_{\frac{n+1}{2}}$, donc le  $\dfrac{n+1}{2}$ entier impair par définition de $i$.
\par Par suite il vient que $n-1$ est un entier pair, il est égal à $p_{\frac{n-1}{2}}=n$, donc le  $\dfrac{n}{2}$ entier pair par définition de $i$.

Ainsi on a les équations suivantes:
$$1+2+3+\cdots+(n-1)+n$$
$$\Longleftrightarrow i_1+ p_1+ i_2 + \cdots + p_{\frac{n-1}{2}} + i_{\frac{n+1}{2}}$$
quitte à réarranger les termes de la somme on a:
$$\Longleftrightarrow (p_1+p_2+\cdots+p_{\frac{n-1}{2}}) + (i_1+ i_2 + \cdots + i_{\frac{n+1}{2}})$$
$$\Longleftrightarrow (\dfrac{n-1}{2}(\dfrac{n-1}{2}+1))+ ( (\dfrac{n+1}{2})^2)$$
d'où 
$$\dfrac{n(n+1)}{2} =\dfrac{n-1}{2}(\dfrac{n-1}{2}+1) +  (\dfrac{n+1}{2})^2 $$

Avec nos suites on peut dire que $$g_{2n+1}=p_{n}+i_{n+1}$$
\newpage

\section{Suites géométriques}
\subsection{Définition}
\begin{defn}
Soit $\suite[u]$ une suite et $q$ un réel.
\par On dit qu'une suite $\suite[u]$ est une suite arithmétique de raison $q$ si pour tout entier naturel $n$ : \[u_{n+1}=q\times u_n\]
\end{defn}

Remarque : Si $q\neq0$ et $u_{0}\neq0$ alors pour tout $n\in\N$, $u_{n}\neq0$.

\begin{exemple}
\par La suite des puissances de 2 est la suite géométrique de premier terme 1 et de raison 2.
\par Les suites $u$ et $v$ définies sur $\N$ par $u_n=-4\times 3^{n}$ et $v_n=n^2+1$ sont-elles géométriques ?


Pour tout $n$, $\dfrac{u_{n+1}}{u_n}=\dfrac{-4\times 3^{n+1}}{-4\times 3^{n}}=3$.
\par La suite $u$ est donc géométrique de raison 3.
\par $v_0=1$, $v_1=2$, $v_2=5$ ainsi $v_1=2\times v_0$ et $v_2=2,5\times v_1$
\par La suite $v$ n'est donc pas géométrique.
\end{exemple}
\subsection{Expression de $u_n$ en fonction de $n$}
\begin{prop}
Soit $\suite[u]$ une suite géométrique de raison $q$.
\par Pour tous entiers naturels $n$ et $p$ : \[u_n=u_p\times q^{n-p}\]
\par En particulier, pour tout entier naturel $n$ : \[u_n=u_0\times q^{n}\]
\end{prop}

\begin{demo}
\par On suppose $n>p$. On a : \\
\par $\dfrac{u_{n}}{u_{n-1}}=q$ ; \hspace{2em} $\dfrac{u_{n-1}}{u_{n-2}}=q$ ; \hspace{2em} \dots \hspace{2em} $\dfrac{u_{p+2}}{u_{p+1}}=q$ ; \hspace{2em} $\dfrac{u_{p+1}}{u_{p}}=q$
\par En multipliant toutes ces égalités, on obtient :
\par $\dfrac{u_{n}}{u_{p}}=q^{n-p}$ soit $u_n=u_p\times q^{n-p}$
\end{demo}

\newpage
\begin{prop}
Soit $\suite[u]$ une suite géométrique de raison $q\neq0$.
\par Si $q>1$ alors 
$\left\{\begin{tabular}{@{}l}
si $u_0>0$, $u$ est strictement croissante.\\
si $u_0<0$, $u$ est strictement décroissante.
\end{tabular}\right.$
\par Si $q=1$ alors $u$ est constante.
\par Si $0<q<1$ alors 
$\left\{\begin{tabular}{@{}l}
si $u_0>0$, $u$ est strictement décroissante.\\
si $u_0<0$, $u$ est strictement croissante.
\end{tabular}\right.$
\par Si $q<0$ alors $u$ n'est pas monotone.
\end{prop}

\begin{demo}
\par Pour tout $n$, $u_{n+1}-u_{n}=u_0\times q^{n+1}-u_0\times q^{n}=u_0\times q^{n}(q-1)$
\par Le signe de $u_{n+1}-u_{n}$ s'obtient donc en fonction du signe de $u_0$, du signe de $q^{n}$ et du signe de $q-1$.
\end{demo}

\subsection{Somme de termes consécutifs}
\begin{prop}
Soit $q$ un réel différent de 1.
\par Pour tout entier naturel $n$ :
\[ 1 + q + q^{2} + \cdots + q^{n-1} + q^{n} = \dfrac{1-q^{n+1}}{1-q} \]
\par Soit $\suite[u]$ une suite géométrique de raison $q\neq 1$.
\par Pour tout entier naturel $n$ :
\[u_{0}+u_{1}+ \cdots + u_{n-1}+u_{n}=u_{0}\times\dfrac{1-q^{n+1}}{1-q} \]
\end{prop}

\begin{demo}
\par Posons $S=1 + q + q^{2} + \cdots + q^{n-1} + q^{n}$. On a alors $qS=q + q^{2} + q^{3} + \cdots + q^{n} + q^{n+1}$
\par Ainsi, $S-qS=1-q^{n+1}$
\par D'où : $S=\dfrac{1-q^{n+1}}{1-q}$
\end{demo}

\begin{exemple}
Déterminer la somme des $n$ premières puissances de 2 ($1+2+2^2+\cdots+2^{n-1}$).


\par La suite des puissances de 2 est la suite géométrique de raison 2 et telle que $u_0=1$
\par Ainsi, $u_{0}+u_{1}+ \cdots+u_{n-1}=u_{0}\times\dfrac{1-q^{n}}{1-q}=1\times\dfrac{1-2^n}{1-2}=2^n-1$

\end{exemple}
\begin{comment}
\newpage

\section{Exercices}
\subparagraph*{Exercice 1}
Soit $u_0=1000$ et  $u_{n+1}=1,04u_n$.
%Alors, $u_0=1000$, \\
%$u_1=1,04\tm u_0=1,04\tm 1000=1040$, \\
%$u_2=1,04\tm u_1=1,04\tm 1040=1081,6$,  \\
%$u_3=1,04\tm u_2= \ \dots$ \\
%$u_{50}=1,04 u_{49}$
Donner $u_0$, $u_1$, $u_2$, $u_3$, $u_{10}$ et $u_{50}$

\subparagraph*{Exercice 2}


Soit la fonction $f$ définie sur $\R$ par $f:x\mapsto 3x-2$. 
\begin{enumerate}
\item On définie la suite $(u_n)$ par $u_0=1$ et  $u_{n+1}=f(u_n)$. 
  Calculer $u_1$, $u_2$, $u_{10}$, $u_{100}$ et $u_{1000}$. 

\item On définie la suite $(v_n)$ par $v_0=2$ et  $v_{n+1}=f(v_n)$. 
  Calculer $v_1$, $v_2$, $v_{10}$, $v_{100}$ et $v_{1000}$. 
\end{enumerate}

\subparagraph*{Exercice 3}
\begin{enumerate}
\item Soit la suite arithmétique $\suite[u]$ de premier terme $u_0 = -5$ et de raison $r = 2$. Calculer $u_{3002}$.
\item Soit la suite arithmétique $\suite[v]$ de premier terme $v_2 = 1200$ et de raison  $r = -10$. Calculer $v_{25}$. À partir de quel rang la suite est-elle négative ?
\end{enumerate}

\subparagraph*{Exercice 4}
Soit $\suite[u]$ une suite arithmétique telle que $u_{10} = -70$ et $u_{25} = 80$.
Calculer la raison $r$ de cette suite, puis calculer $u_0$ et$u_{1212}$.

\subparagraph*{Exercice 5} 
Soit  $\suite[v]$ la suite géométrique de premier terme $u_0 = 0,2$ et de raison $q =\dfrac14$.
Calculer $u_4$ et $u_{20}$.
\subparagraph*{Exercice 6} On utilise une feuille de papier, d'épaisseur $e = 0, 5 mm$, que l’on replie successivement en deux. Quelle est l’épaisseur de la feuille après le premier pliage ? après le deuxième ? après le $n$-ième ? Combien de fois faudrait-il replier cette feuille en deux pour obtenir une épaisseur supérieure à $300 m$ ?

\end{comment}
\end{document}
